% ============================================================
% 09_conclusiones.tex — Conclusiones y Trabajos Futuros
% ============================================================

\section{Conclusiones}
\label{sec:conclusiones}

El presente trabajo presentó el diseño, implementación y validación de NodeAlert IoT,
un sistema embebido de monitoreo ambiental y control autónomo de iluminación para
entornos educativos. Los resultados obtenidos permiten enunciar las siguientes
conclusiones:

\begin{enumerate}
  \item \textbf{Viabilidad del enfoque edge computing en control de iluminación}:
        la implementación de la lógica de control completamente en el firmware del
        microcontrolador demostró ser una elección correcta. Durante las pruebas,
        el sistema mantuvo el control de la iluminación de forma autónoma incluso
        en ausencia de conectividad MQTT, confirmando la independencia funcional
        de la capa de percepción respecto de la capa de aplicación.

  \item \textbf{Efectividad de la lógica de histéresis de doble umbral}:
        la función \texttt{regla\_deseado()} demostró comportamiento correcto en
        todos los escenarios de prueba definidos, incluyendo los casos de falla
        de uno o ambos sensores de luz, la gestión del lazo de retroalimentación
        inherente al sensor de aula y la transición suave mediada por el debounce
        de 2\,s.

  \item \textbf{Integración exitosa del radar de microondas LD2420}:
        el driver desarrollado en C puro sobre ESP-IDF, incluyendo el parser de
        tramas nativas y el detector de movimiento por votación, operó
        establemente durante las pruebas. La selección del radar frente a sensores
        PIR convencionales aporta una mejora cualitativa en la detección de
        personas estáticas relevante para el contexto de aulas universitarias.

  \item \textbf{Arquitectura de servidor escalable y reproducible}:
        el stack Docker Compose (Mosquitto + Node-RED + React/Nginx) permitió
        desplegar toda la infraestructura del servidor con un único comando,
        con separación clara de responsabilidades entre servicios. La generación
        programática de flujos Node-RED mediante \texttt{generate-flows.js}
        garantiza la reproducibilidad del despliegue.

  \item \textbf{Cumplimiento de los requisitos de la cátedra}:
        el sistema satisface la totalidad de los requisitos mínimos establecidos
        por la materia Integración Computacional I e II: programación de
        microcontrolador con optimización de recursos y ahorro de energía,
        comunicación bidireccional por MQTT con broker propio configurado,
        suscriptores y publicadores diferenciados, e interfaz de usuario con
        visualización de variables y controles remotos.
\end{enumerate}

\section{Trabajos Futuros}

A partir de los resultados obtenidos y las limitaciones identificadas, se proponen
las siguientes líneas de trabajo futuro:

\begin{enumerate}
  \item \textbf{Configuración dinámica de umbrales}: implementar la modificación
        de los parámetros de control (umbrales de luminosidad y tiempo de retardo)
        en tiempo de ejecución, almacenando los valores en NVS Flash y
        actualizándolos mediante un mensaje MQTT específico desde el servidor,
        eliminando la necesidad de reflashear el firmware.

  \item \textbf{Comunicación TLS/SSL}: agregar soporte para MQTT sobre TLS en el
        cliente ESP-IDF y en Mosquitto, con gestión de certificados integrada en el
        contenedor Docker, para garantizar confidencialidad e integridad de la
        telemetría en redes no confiables.

  \item \textbf{Validación multi-nodo}: extender el prototipo a múltiples nodos
        físicos distribuidos en diferentes aulas, validando la escalabilidad de la
        arquitectura del broker y de la base de datos SQLite.

  \item \textbf{Calibración de sensores}: comparar las lecturas de los sensores
        BH1750 y SHT30 contra instrumentos de referencia calibrados (luxómetro y
        termohigrómetro certificado) y aplicar factores de corrección si fuera
        necesario.

  \item \textbf{Visualización ampliada}: agregar gráficos de serie temporal para
        iluminancia (ventana y aula) y estado del relé en la vista NodeDetail, y
        exportación de datos históricos en formato CSV.

  \item \textbf{Integración con sistemas de gestión de edificios (BMS)}: definir
        una interfaz de integración con plataformas estándar de gestión de
        edificios (como Modbus TCP o BACnet/IP) para facilitar la incorporación del
        sistema en infraestructuras de edificios inteligentes existentes.

  \item \textbf{Ahorro energético cuantificado}: instrumentar el consumo eléctrico
        de la carga de iluminación controlada durante un período de tiempo
        representativo y compararlo con una instalación sin control automático, para
        cuantificar el ahorro energético real del sistema.
\end{enumerate}

\section*{Agradecimientos}

Se agradece a los docentes de la cátedra de Ingeniería en Computación II de la
Universidad Nacional de Rafaela (UNRaf), Milton Pozzo y Matías Wanzenried,
por su orientación técnica, seguimiento y aportes metodológicos durante el
desarrollo de este proyecto.

